\section{Existing Acoustics Free Open Source Software}

Before discussing the development of the \texsc{magpie} framework it is worthwhile discussing the state of free open source software in the field of acoustics.
Section~\ref{sec:aurora} has already covered the case of the Aurora plug-ins created by Angelo Farina and the influence of that software of the open methodology of \textsc{magpie}.
This section will look at a broader set of acoustics software, all of which would fall under the umbrella of FOSS, though maybe not explicitly.
Taken as whole they should provide an insight into FOSS practice within acoustics and what areas the community needs to develop.

An overview of specifically open source room acoustics software can be found in \textcite{hornikx_exploring_2024}.

Software discussed here has direct application in the heritage sector, either through measurement of heritage architecture \cite{tronchin_acoustics_1997,fagerlande_adequaccao_2020} or through digital reconstruction \cite{katz_exploring_2020}


Measurement provided by tools such as the ITA toolbox, aurora. Simualation provided by NeSS, PFFDTD 

``Exploring cultural heritage through acoustic digital reconstruction provides historians, musicologists, and others with a perspective not available using more established research methods. Furthermore, it brings a powerful means of communicating and delivering memorable, meaningful, and most importantly, multisensory experiences. The effectiveness of digital reconstruction is evident through the range of projects undertaken throughout Europe'' \cite{katz_exploring_2020}

Each example may not have a direct application within the heritage sector, but as is shown with Aurora, there is no reason that they could not be.

First a clarification, when discussing documentation it is not just the instructions for the use of the software but also the way that is structured. A interconnected set of documents that describe the class hierarchies and function usage.
The documentation is ideally generated from the source code though documentation comments, this stops duplication, but also provides a means for automation as comments.
There are many documentation standards in one single language let alone across the three programming languages discussed in this chapter: MATLAB, Python, and C\textsc{++}. 
MATLAB has its style for comments and generating documentation from them while Python \cite{vanrossum_pep8_2001,goodger_pep257_2001} and C++ have a number of competing standards
\footnote{Doxygen, Sphinx, HeaderDoc to name a few documentation systems with respective documentation syntax. Distribution of popularity for these documentation systems would be difficult to ascertain without a comprehensive analysis of open source projects. The popular computer programming question-and-answer site stackoverflow provides open meta-data exploration and a query of tags reveals Doxygen to be a clear frontrunner followed by Sphynx.\\\url{https://data.stackexchange.com/stackoverflow/query/1069131/get-all-tags}}
There is a stress \textit{for} documentation for as part of best practice in software sustainability \cite{smith_software_2016}, but less on what that documentation should contain. 

from industry \cite{google_documentation_2025}

Generating documentation comments can that can use markup like HTML to enable cross referencing and easier searching of the source that can be done separate to the development environment

The documentation does not just serve an outer community but also provides an aid memoire for
``Understand your code in 6 months'' \cite{writethedocs_sowftware_2025} 


\textcite{writethedocs_sowftware_2025} provides a comprehensive guide on documentation best practice

Documentation is aimed at two audiences, users and developers who can sometimes be the same person.

Documentation should be there to aid developers and user alike, who might sometimes even be the same person.

There is no strict guideline on exactly what developer documentation should contain. A general rule being that documentation should provide a means to ``understand your code in 6 months'' \cite{writethedocs_sowftware_2025} with the consequences of not adhering this rule illustrated multiple times in \textcite{perkel_challenge_2020}.

Such developer documentation is typically contained directly in the comments of the source code. Generating documentation comments can that can use markup like HTML to enable cross referencing and easier searching of the source that can be done separate to the development environment.


Documentation comments strings are not manual, but rather references 

\url{https://www.software.ac.uk/blog/code-works-isnt-always-code-lasts}
\cite{procida_diataxis_2025} 4-point framework of Reference, How-To, Tutorials, Explanation.

\cite[Section 6.2]{stallman_gnu_1992} warns against the conflation of developer ``Doc Strings'' and a user manual. A stand alone user manual \cite{lehoucq_arpack_1998, },


DOI for the software as a whole and one for each release \cite{smith_software_2016}


https://www.software.ac.uk/guide/good-practice-submitting-software-outputs-ref2021
source code;
compiled code and/or instructions on how to compile or install the software;
user documentation;
test suite and/or quality assurance procedures.

DOI for a companion software paper;
DOIs for research outputs resulting from use of the software;
Usage metrics.

Requirements

\url{https://spdx.org/licenses/} SPDX License List

Some requirements not explicitly stated, with a version control repository there also typically comes issue tracking and contribution functionality. Similarly, if a version control repository is absent then a project is also unlikely to contribution and code of conduct.

\begin{table}[]
    \centering
    \begin{tabular}{c|c}
         &  \\
         & 
    \end{tabular}
    \caption{Acoustics FOSS compliance Overview}
    \label{tab:placeholder}
\end{table}


```Some programming systems, such as Emacs, provide a documentation string for each func- tion, command or variable. You may be tempted to write a reference manual by compiling the documentation strings and writing a little additional text to go around them—but you must not do it. That approach is a fundamental mistake. The text of well-written documentation strings will be entirely wrong for a manual.
A documentation string needs to stand alone—when it appears on the screen, there will be no other text to introduce or explain it. Meanwhile, it can be rather informal in style.
The text describing a function or variable in a manual must not stand alone; it appears in the context of a section or subsection. Other text at the beginning of the section should explain some of the concepts, and should often make some general points that apply to several functions or variables. The previous descriptions of functions and variables in the section will also have given information about the topic. A description written to stand alone would repeat some of that information; this redundancy looks bad. Meanwhile, the informality that is acceptable in a documentation string is totally unacceptable in a manual.
The only good way to use documentation strings in writing a good manual is to use them as a source of information for writing good text.``

- Reference (Developer) Documentation
- A User Guide
- A Unique Identifier (DOI)
- A Version controlled

Documentation should be comprehensive to multiple modes, for usage, for development, and should be also proved functioning self-contained examples \cite{aghajani_software_2020}.
references what should be in documentation for hardware \cite{oshwa_best_2013}


Creative Commons discourages usage of its licence for software \cite{cc_can_2025}
\footnote{Although CC-BY-NC-SA-4.0 is ``one-way'' and CC0 fully GPL ``compatible'' they are still discouraged by the Creative Commons \cite{cc_faq_2016} } and no creative commons licence is Open Source Initiative approved \cite{osi_licenses_2025,spdx_licenses_2025}.

https://spdx.org/licenses/

Taking all that has been discussed so far and using the \citp{ssi_checklist_2018} checklist we can narrow down the core requirements of an open source research project as source code availability, a version control system (VCS), documentation (developer reference and user manual, a compliant open source licence and a unique identifier, typically a DOI. There is a distinction here between the source availability (the ``open'' in ``open source'') and the presence of a version control system. A FOSS project can consist of all the base elements listed above but may not choose to have a public version control system. A pubic VCS by definition implies an availability of the source, but the reverse is not automatically true, which we will see with the projects discussed in this section. Similarly a unique identifier such as a DOI implies digital archiving has taken place but, again, not necessarily the reverse.

Recommendation for research software extend beyond these categories, but they are dependant on their existence. For instance, contribution policies, issue tracking and code of conduct only apply if the source code is available and VCS has been implemented. Open source projects should first be the base requirements before they are assessed for more comprehensive compliance. 

This section will discuss seven open source research software projects pertaining to the field of acoustics, of which applications within musical acoustics and further have viable usage within musical cultural heritage studies.


Examples are:

\item VK-Gong: \url{https://vkgong.ensta-paris.fr}
\item ITEMM Open Lab no repo, no doi, no licence, no documentation
\item NESS: \url{https://github.com/Edinburgh-Acoustics-and-Audio-Group/ness} open source git repository of the Next Generation Sound Synthesis project. A suite of software for the modelling musical instrument systems and acoustic spaces. No user manual, no DOI, version controlled, compiler instructions for arm architecture
\item PFFDTD

Vk-gong \url{https://vkgong.ensta-paris.fr} extensive documentation \url{https://vkgong.ensta-paris.fr/files/VKgong1.0_documentation.pdf} based on research in \cite{ducceschi_simulations_2015}
\cite{nguyen_nonlinear_2019} ``The input parameters of the force are thus the time t0, the temporal width of the interaction Twid, and the amplitude of the force pm in N. All the computations have been implemented using the software vk-gong,25 an open source code developed to handle nonlinear vibrations of plates and built from previous works by the second author.''

and \cite{diaz_fast_2025}

``The open source repository includes utilities for computing analytical modes and eigenvalues of fixed strings, membranes and plates. For plates, we provide a Python implementation of the approach described in [11] to compute non-linear coupling coefficients``

\subsection{Aurora}

The status of Aurora has already been discussed Section~\ref{sec:aurora} in the context of its influence on the open methodology applied to the projects of this thesis. In terms of meeting the base components, compliance has only happened relatively recently. The source is open and
, \cite{farina_new_2009} Created in 1995, Aurora is the oldest FOSS project among the sample set. Although ``free''
\footnote{as in beer, not as in freedom.}
the ``open'' only began to start in 2009 with the porting of Aurora to the Audacity open source audio editor \cite{farina_new_2009}. The source was open on request and would be shared by the authors, but was not available for direct public download.

With the transition to an entirely open source code base \cite{farina_aurora_2025}, Aurora now has a unique identifier for citation, a open source initiative compliant licence. Documentation is lacking, however, with the source code being largely under documented and no developer reference that can be consulted. Similarly, a user manual exists of sorts, but it is not open to contribution and exists, currently, on Farina's Aurora Plugins website,
\footnote{\url{https://www.aurora-plugins.com/Package.htm}}
the future of which is uncertain.


\subsection{ITA Toolbox}

The ITA toolbox  had its genesis in 2010 \cite{dietrich_matlab_2010} and migrated to the RWTH Aachen hosted GitLab VCS in 2016 and announcement of  an explicit open source model in 2017 \cite{berzborn_ita_2017}. ITA Toolbox can be viewed as a successor to the aurora plug-ins \cite{dietrich_influence_2013}, providing tools for impulse response measurement and analysis. 


Source code developer documentation is extensive and uses MATLAB's internal documentation system and standards but, as warned by \textcite{stallman_gnu_1992}, ``doc strings'' are not a substitute for a user manual. The user documentation is limited to what would be deemed ``Tutorials'' in the \citetitle{procida_diataxis_2025} \cite{procida_diataxis_2025}, but the other three documentation forms of ``How-To``, ``Explanation`` and ``Reference`` are more limited. What is primarily missing from the ITA Toolbox is documentation that exists outside of the M\textsc{ATLAB} environment and more openly available for consultation. 

The user manual and citation system are entertwined as part its publication historu
Recommended citation is through ITA Toolboxes publication history
Source exists under a single, release https://git.rwth-aachen.de/ita/toolbox/-/releases

Application in assessment of heritage architecture in \textcite{fagerlande_adequaccao_2020}, which discussed the architectural adaption of the heritage building Teatro Armando Gonzaga in Rio di Janeiro through simulation for which the impulse responses were facilitate by the ITA Toolbox.  

\textcite{berzborn_ita_2017} states the project's intention to register a DOI via the Zenondo platform, but as of writing this has not yet been carried out.

It is worth noting also that operates with the proprietary MATLAB environment and requires the additional DSP System, Signal Processing  and Curve Fitting Toolboxes restricting users to only those whose licences cover these additional dependencies. Although the source is open and can be ported to any language, the use of proprietary tool boxes makes this a much less practical undertaking. On a final note, the variant of the BSD-4-Clause License
\footnote{\url{https://spdx.org/licenses/BSD-4-Clause.html}}
used for ITA Toolbox is not currently approved by the open source initiative. 


ITA Toolbox represents one of the most successful open source endeavours in this sample set and is the only project whose version control is hosted by its institution.



\subsection{ITEMM OpenLab}

% \cite{viala_guide_2021,charrier_lopen_2023}

% ``pour la filière des instruments de musique.''\cite{charrier_lopen_2023}

% ``Tools:

% This part will give access to online tools that measure the rigidity of plates, beams, wooden quarters, more generally to identify the properties of materials. It will also offer a module for piano strings and the piano string plan. This section will allow you to calculate the properties of composite materials in advance.

% Finally, an interface will offer the ability to extract the audio characteristics of a recorded sound.''

In its legal text, OpenLab states that its  work under a CC-BY-SA though the version is not explicitly stated this is an implied CC-BY-SA-4.0.


\subsection{NeSS}
The Next Generation Sound Synthesis project from the University of Edinburgh is a suite of physically modelled sound synthesis algorithms. For these algorithms, Graphics Processing Units were leveraged for large scale, parallelised numerical simulations such as brass instrument, string / fretboard interaction and virtual room acoustics. NESS was a joint project between the Acoustics and Audio Group at the University of Edinburgh and the Edinburgh Parallel Computing Centre, running for five years between ran between 2012 and 2017. 

During the course of this thesis the source code for NeSS was made publicly available under an OSI approved MIT Licence via the online VCS service GitHub.


Like ITA Toolbox, User manual exists in some for through its publication http://www.ness.music.ed.ac.uk/music-and-tools/publications


Cultural heritage usage \cite{katz_exploring_2020}


NeSS has implemented alternative compilation instructions for newer arm processors.

\cite{bilbao_physical_2019}

\subsection{PFFDTD}
\cite{hamilton_pffdtd_2021}
provides: 

- suggested citation
- background references

\subsection{VK-Gong}

VK-Gong is software created out of ENSTA Paris for the numerical simulation of nonlinear plates and shells using Föppl–von Kármán theory for nonlinear vibrations. Such non linear systems would be found in percussive instruments such as gongs, cymbals and bells. 

VK-Gongs provides its source code, in MATLAB and C\textsc{++}, from an ENSTA Paris hosted site. The project hosts one of the most comprehensive user manuals \cite{vkgong_manual_2017} as part of this sample set, which provides the technical background and usage instructions in both a `how-to' and tutorial format, although this applies entirely to the MATLAB version of the source. The C\textsc{++} is marked as ``preliminary'' with documentation to be ``ready in summer 2017.''

VK-Gong's site states the work under is under the CC-BY-NC-SA-4.0 licence,
\footnote{\url{https://creativecommons.org/licenses/by-nc-sa/4.0/}}
although the license text does not appear within the source as either a document string or a separate license file. The usage of CC licenses for software is discouraged by both the Creative Commons and Open Source Initiative and though the CC-BY-SA-4.0 is ``one way'' compliant with the GPLv3 \cite{cc_can_2025}, the CC-BY-NC-SA-4.0 is not.

Though the source is openly available from the VK-Gong site, it does not use any version control system. Without a VCS, the growth of the software is stunted which is a missed opportunity especially given the C\textsc{++} version of VK-Gong is still in its preliminary stages. The CC-BY-NC-SA-4.0 licence permits the mirroring of the source code,
\footnote{\url{https://github.com/rodrigodzf/VKGong}}
uploading of this code onto an online archive platform like Zeonodo is not compliant with its rules \cite{zenodo_metadata_2025}. 

The main technical obstacle with VK-Gong source code is the usage of SSE intrinsic instructions, which are only applicable to intel processors. The SSE instructions NeSS source code, separate compilation instructions have been provided for specific CPU architecture and operating system.



\subsection{YIN}
The YIN pitch detection algorithm was also briefly discussed in Section~\ref{sec:aurora} alongside the case study of Aurora as example open source acoustics software practice.

The theory behind the YIN algorithm is discussed in \textcite{cheveigne_yin_2002} and its MATLAB source code code see direct usage in research.

Like Aurora, YINs source code has no single fixed URI, let alone a DOI with which it can be cited and traced. YIN also also has one of the most restrictive licences with the original MATLAB source limited to ``FOR RESEARCH PURPOSES''
\footnote{\url{https://github.com/mhamilt/yin/blob/main/LICENCE}}
which is not an OSI compliant clause and would disqualify its categorisation as FOSS.

Used in \textcite{bowen_assessing_2019} as part of a study on musical instrument heritage clarinets that were not playable.
\footnote{\textcite{bowen_assessing_2019} cites \textcite{cheveigne_yin_2002} but also remarks that the MATLAB was used as part of the study, making highly likely that the paper used the original YIN source code.}


The documentation for YIN is in the form of the associated publication \cite{cheveigne_yin_2002}, which would constitute the ``Explanation'' form of documentation in \textcite{procida_diataxis_2025}, developer documentation is absent.
Like the problematic SSE intrinsics in VK-Gong, YINs technical stumbling block is the usage of MATLAB's MEX compilation system \cite{matlab_mex_2025}. In YIN's currently form, as it was 2003, the compiled \texttt{.mex} files will not execute correctly in any MATLAB version from R2016a onwards. Instead, the \texttt{.mex} must be recreated due to changes in the compilation system. Recompiling does not provide a great obstacle, and the correct syntax can be quickly found by consulting the documentation, but this becomes a problem that \textit{every} user needs to solve. Not all researcher will have the same technical background and skill-set and without a centralised source, like a version controlled repository, there is no single place to raise or document such technical problems.


\subsection{Summary of Acoustics Software Practice}


From the set of seven acoustics software we can we can see a inclination towards an open source model, but not yet full compliance with current guidance.

Unique identifiers are absent across all examples, though some such as \citetitle{hamilton_pffdtd_2021}, \citetitle{berzborn_ita_2017} have taken step towards citablility by providing an explicit citation.

This is a step in the right direction, but it makes an assumption that software will stay static. The dual DOI for version and for the base software provides the ability to discuss variations of the software. Software functionality may grow or become more accurate or conversely, bugs within a programme may be present in on version and it should be easier to find if research was carried out on potentially malfunctioning software.
Lack of unique identifers also identified in \textcite{hornikx_exploring_2024}.


When citations are distributed across multiple papers such as for the ITA Toolbox, there is unnecessary complexity added to the process of tracking usage and cross referencing, but its is still achievable.

For those without a suggested citation, tracking and cross referencing tends to settle on a finite set of urls or specific articles such as for Aurora and YIN. The usage is difficult to differentiate as a without further scrutiny as these may be citations to the article's contents and not necessarily the usage of the software.

Licensing is fairly comprehensive, though there is some misapplication of Creative Commons licenses where \cite{cc_can_2025,hong_choosing_2025}. Though the BSD 4-Clause is not OSI compliant 
\footnote{\url{https://www.gnu.org/licenses/bsd.en.html}}


As the methodology applied to the keyboard interfaces provides an example of open hardware practice for a research product, so to is \textsc{magpie} intended to serve as example for open source research software.






